% Français 9115, batch 12 — Gallica views 221–240.
% Pass: Opus 5 (claude-opus-5), 2026-09-19 — first pass, unchecked against
% the views by a human.
%
% What the twenty views hold. The Latin fair copy of Euler's « Investigatio
% motuum quibus laminæ et virgæ elasticæ contremiscunt », which filled
% batches 8 to 10, has stopped before this batch begins. What these views
% carry is a piece in French, in the writer's own fair hand: first the close
% of a long quotation — a French version of the end of Euler's memoir, §§ 47
% to 51, set out line by line with a guillemet repeated at the head of every
% line, as she does elsewhere in the volume — and then, from view 228, her
% own critical discussion of it, in numbered remarks. The discussion argues
% that Euler was wrong to admit the second solution of his problem of the rod
% pinned at an intermediate point: the condition sin ω = 0, which he had
% established must hold for every possible motion, cannot be reconciled with
% the second factor of his final equation except for ω = 0, and she works the
% general case λ = 1/λ' out to show it. The batch ends in the middle of that
% argument, at the foot of her page 11.
% The first view of the batch, 221, is a blank verso, and its recto — the
% leaf inked 107 — is view 220, in batch 11; the quotation therefore runs on
% from that batch into view 222 without a break, and the writer's pagination
% of this piece begins there too (her pages 1 and 2 are not in this batch).
%
%   v. 222–226  The end of the quoted French version of Euler: §§ 47–48, the
%               eight equations I–VIII for the two portions and their
%               reduction to one transcendental equation (v. 222–224); the
%               title « Développement du cas dans lequel la lame E F est
%               appuiée non seulement a ses extrémités E et F, mais encor
%               dans le milieu de dela longueur » and § 48's factorisation
%               (v. 224); § 49, the first factor, sin ½ω = 0, giving case IV
%               for each half and the sounds 4, 16, 36, 64 (v. 224–226);
%               § 50, the coefficients, all vanishing but γ = −γ' (v. 226);
%               § 51, the second factor, tang ½ω = (e^ω−1)/(1+e^ω), the same
%               equation as case V, and the sounds 5π/2, 9π/2, 13π/2 (v. 226).
%   v. 228      The quotation closes three lines in, on a closing guillemet,
%               and the writer's own text begins, marked « I. »: the second
%               solution « sort de l'état de la question », since the
%               condition sin ω = 0 that Euler established for every motion
%               agrees with sin ½ω = 0 but not with tang ½ω = (e^ω−1)/(e^ω+1)
%               unless ω = 0. « Cependant le respect dû … l'opinion d'Euler
%               m'a engagé a examiner le cas general. »
%   v. 230      Remark « 2. »: the general equation divided by e^{2λω} − 1;
%               the two classical series for sin ω in powers of sin λω,
%               according as the denominator of λ is odd or even; sin λω is a
%               factor, ω = iπ/λ, and the lamina divides into 1/λ parts that
%               vibrate like so many separate laminæ, the points of rest
%               « que l'on pourroit nommer sympatiques ».
%   v. 232      A larger leaf mounted on a guard, inked 113, with no page
%               number of the writer's, written only on the top fifth: a
%               first draft of the two opening paragraphs of her page 8
%               (view 234), the coefficients α, β, γ, δ and α', β', γ', δ' in
%               the general case. Its first line does not follow the last of
%               view 230.
%   v. 234      Page 8: the sounds of the fractional parts, πc√2gb/λλaa, 4, 9,
%               16 times it; the same coefficients, better set out; and the
%               equations of the successive parts E L, L E', L' L'', L'' L''',
%               with the sign + or − according as the part lies above or below
%               the line of rest. A marginal note sends the reader to fig. 7.
%   v. 236      Page 9, remark « 4. »: the second factor of the equation of
%               the problem written out with the two series, and the proof
%               that it cannot vanish, when sin ω = 0 is respected, except for
%               ω = 0. Two marginal notes.
%   v. 238      Page 10, remark « 5 »: the final equation transformed in four
%               steps into sin λω{…} + sin(1−λ)ω{…} = 0, and the two
%               conditions tang(1−λ)ω = …, tang λω = … that would have to hold
%               together if the pin could be treated as a fixed point — which
%               is what Euler assumes in his second solution, says a marginal
%               note.
%   v. 240      Page 11: the two series of sounds for the portions of length
%               λa and (1−λ)a, the arithmetic of λ = 1/λ' worked out to
%               λ' = 6, and remark « 6 », the conclusion: « J'ai voulu entrer
%               dans tous ces details afin d'être autorisée a conclure
%               qu'Euler a été induit en erreur lorsqu'il a cru pouvoir
%               admettre sa seconde solution… ». The page is the most heavily
%               corrected of the batch. The text breaks off here.
%
% Views not transcribed. Ten: 221, 223, 225, 227, 229, 231, 233, 235, 237 and
% 239. Every one is the blank verso of the leaf photographed on the view
% before it; what looks like writing on them is the recto showing through,
% and reads reversed at full resolution — the ink number of the recto shows
% through at the top left of each. No view of this batch is a second
% exposure, and none is a microfilm target.
%
% The numbers on the leaves, and why no \folio{} is written. Two kinds, both
% in ink, one to a written leaf. (1) The writer's own pagination of this
% piece, at the middle of the top of the page and underlined: 4 at view 224,
% 5 at 226, 6 at 228, 7 at 230, 8 at 234, 9 at 236, 10 at 238, 11 at 240. At
% view 222 the figure is cut by the folded edge of the leaf and is read 3
% with doubt. View 232 carries none. (2) The ink series of the earlier
% batches, at the top right of the written side: 108 at view 222, 109 at 224,
% 110 at 226, 111 at 228, 112 at 230, 113 at 232, 114 at 234, 115 at 236, 116
% at 238, 117 at 240 — one number per leaf, without a gap. No pencilled
% foliation of the Bibliothèque was read anywhere in the batch, so no
% \folio{} is written, as in batches 1 to 10. Every number above was read on
% its view; none is inferred.
% A third series is on the leaves but is not a number of the leaves: the
% writer numbers her remarks in the margin or at the head of the line —
% « I. » at view 228, « 2. » at 230, « 2 » at 232 and at 234, « 4. » at 236,
% « 5 » at 238, « 6 » at 240. The three marks read 2 are the same figure,
% stroke for stroke, and none of them is a 3; the run is transcribed as it
% stands and not made consecutive.
%
% Notation. The letters are the copy's: λ for the fraction EL/EF, ω, u, ζ,
% C, α β γ δ and their primes, a, b, c, g. The radical is written without a
% bar and is set \surd. « sin. », « cos. », « tang. », « cot. » are written
% with and without their stop in the same line, and the pass has not tried to
% record every stop. The looped sign closing a series is set « \&c. ». The
% writer's run-together words are kept: « dela », « lamoitié », « lamême »,
% « alaquelle », « lalongueur », « delames », « aune », « laquestion »,
% « nepeuvent », « audessus », « audessous », « aumoins », « cidessus »,
% « apresent », « deplus », « nonseulement ». So is her spelling:
% « coëfficiens », « mouvemens », « appuiée », « stilet », « sympatiques »,
% « deffinitive », « parceque », « adire », « pur » for « pour ». Square
% brackets in the text are hers. Where the calculation and the leaf disagree
% the leaf stands, with a note: view 222 (γ without its denominator, a « − »
% where the memoir has « + », the second end written E where the memoir
% writes F), view 224 (a verb missing after « réduite »), view 226
% (« s = u/a »), view 232 (cos ω = a blot), view 236 (2/λ against 1/λ).
%
% What could not be read. \ill{} is written at: view 224, the whole formula
% struck out with a heavy zigzag after « réduite a celle ci »; view 228, the
% letter opening the small parenthesis « (… cas IV §. 35 ) » and the word
% struck after « établi » and after « dû au »; view 230, the word struck
% after « solution » and the words struck after « examiné »; view 232,
% nothing; view 238, the words struck after
% « en même tems », after « auroit alors », the end of the page, and the
% word after « que cette portion »; view 240, the head of the numerator of
% (1−λ)², two whole struck lines read only in part, and the word struck in
% « j'ai … voire ». Readings offered with doubt are marked \uncertain{}: the
% page number at view 222; the paragraph number 49 at view 224; « 2 » over λ
% at view 236; « cos. », « donc », « seroit », « pour » at view 238; « 1 »
% twice, « supposé », « voire », « dernière » at view 240; « 1 » at view 232.
% Nothing was guessed to fill a gap.
%
% Nothing here was checked against any published transcription or edition of
% Euler's memoir or of Sophie Germain's papers.
\documentclass[11pt,a4paper]{article}
% The preamble shared by every transcription of the mathematicians' manuscripts in Gallica.
%
% It rests on one idea: **uncertainty compounds, it does not resolve.** A
% transcription that smooths over a doubtful word has destroyed the only thing
% it was made for — a reader must be able to tell, at every point, what was on
% the page and what was supplied. The apparatus macros below are therefore the
% critical apparatus, not decoration.
%
% The same commands are recognised by `scripts/render.mjs`, which produces the
% site's reading view, and by `scripts/tei.mjs`. A macro added here must be
% added there too, or rendering fails loudly — which is the wanted behaviour.
%
% Comments here are English, like the rest of the repository. The documents
% this preamble sets are French, and that is deliberate: see the skills.

\ProvidesPackage{germain}[2026/09/15 Transcriptions of mathematicians' manuscripts in Gallica]

\usepackage{amsmath,amssymb,amsthm}
\usepackage{mathrsfs}
% Kept from the parent preamble so the renderer's diagram support stays
% available; these manuscripts draw no commutative diagrams, and nothing here uses it.
\usepackage{tikz-cd}
\usepackage[english,french]{babel}
\usepackage{fontspec}

% --- Greek ------------------------------------------------------------------
% The text face stays fontspec's default, Latin Modern, which has no Greek:
% XeTeX drops every glyph it lacks with a line in the log and nothing on the
% page, and the Greek words the manuscripts quote vanished from the PDFs. So
% the Greek and Greek Extended blocks get a XeTeX character class of their own,
% and entering or leaving a run of it switches the family to CMU Serif — the
% same Computer Modern design, with polytonic Greek — and back. Only the
% family changes: a Greek word in \textit or \textbf keeps its shape and series.
% The transitions to and from every other class carry the tokens that class
% has towards an ordinary letter, so French spacing before « ; » or inside
% « » still holds after a Greek word. No grouping, so no brace or \endgroup
% can come out unbalanced; maths is left alone, since XeTeX inserts nothing
% there. Kept identical in `exercices.sty`.
\makeatletter
\newfontfamily\germainGreekfont{cmun}[
  Extension=.otf, NFSSFamily=germaingreek,
  UprightFont=*rm, ItalicFont=*ti, SlantedFont=*sl,
  BoldFont=*bx, BoldItalicFont=*bi, BoldSlantedFont=*bl]
\newXeTeXintercharclass\germain@greek
\def\germain@greekrange#1#2{%
  \@tempcnta=#1\relax
  \loop
    \XeTeXcharclass\@tempcnta=\germain@greek
    \ifnum\@tempcnta<#2\advance\@tempcnta\@ne\repeat}
\germain@greekrange{"0370}{"03FF}
\germain@greekrange{"1F00}{"1FFF}
\def\germain@greekprev{\rmdefault}
\def\germain@greekin{%
  \edef\germain@greekprev{\f@family}\fontfamily{germaingreek}\selectfont}
\def\germain@greekout{\fontfamily{\germain@greekprev}\selectfont}
\def\germain@greekjoin{%
  \XeTeXinterchartokenstate=\@ne
  \@tempcnta=\z@
  \loop
    \ifnum\@tempcnta=\germain@greek\else
      \XeTeXinterchartoks\germain@greek\@tempcnta=\expandafter{%
        \expandafter\germain@greekout\the\XeTeXinterchartoks\z@\@tempcnta}%
      \XeTeXinterchartoks\@tempcnta\germain@greek=\expandafter{%
        \the\XeTeXinterchartoks\@tempcnta\z@\germain@greekin}%
    \fi
    \ifnum\@tempcnta<4095\advance\@tempcnta\@ne\repeat}
% babel-french sets its own classes, and saves and restores the interchar
% state, each time a language is selected: join again after it does.
\addto\extrasfrench{\germain@greekjoin}
\addto\extrasenglish{\germain@greekjoin}
\AtBeginDocument{\germain@greekjoin}
\makeatother

\usepackage{geometry}
\geometry{a4paper,margin=25mm,marginparwidth=28mm}
\usepackage{marginnote}
\usepackage{xcolor}
\usepackage[normalem]{ulem}
\setlength{\emergencystretch}{3em}
\usepackage{hyperref}
\hypersetup{colorlinks=true,linkcolor=black,urlcolor=black,pdfborder={0 0 0}}

% --- Batch metadata ---------------------------------------------------------
% Taken as they stand by the rendering script, which shows them at the head of
% the reading view. They fix what the file claims to cover. The macro names are
% the parent project's, so that the scripts stay shared; \folder here means the
% volume's slug — `fr-9115` — and \pages counts Gallica views.

\newcommand{\folder}[1]{\gdef\grothFolder{#1}}
\newcommand{\batch}[1]{\gdef\grothBatch{#1}}
\newcommand{\pages}[2]{\gdef\grothFirst{#1}\gdef\grothLast{#2}}
\newcommand{\foldertitle}[1]{\gdef\grothFoldertitle{#1}}

% The holder's own shelfmark, printed as they write it: « Français 9115 ».
\newcommand{\shelfmark}[1]{\gdef\grothShelfmark{#1}}

% The dating, copied verbatim from the holder's catalogue — never inferred
% here. « XIXe siècle » is the BnF's; a pass that dates a leaf from its content
% says so in a \note{}, as its own inference.
\newcommand{\dating}[1]{\gdef\grothDating{#1}}

% The status notice, printed in the title block and repeated as a diagonal
% watermark on every page of the PDF. The manuscripts are in the public
% domain; what this line declares is not a right but a quality — an
% unverified first pass by a machine. The skills write the line; a file
% without it gets no watermark, so leaving it out is visible in the output.
\newcommand{\watermark}[1]{\gdef\grothWatermark{#1}}

\gdef\grothFolder{?}\gdef\grothBatch{}\gdef\grothFirst{?}\gdef\grothLast{?}%
\gdef\grothFoldertitle{}\gdef\grothShelfmark{}\gdef\grothDating{}\gdef\grothWatermark{}

% --- The résumé -------------------------------------------------------------
\newenvironment{resume}
  {\small\setlength{\parskip}{0.45em}}
  {\par\medskip\hrule\medskip}

% --- Keywords ---------------------------------------------------------------
% In English, closing the modernised reading's résumé: the modern vocabulary
% under which to search for what the leaves construct. The manifest extracts
% them to tag the volume — this line is the single source of the tags.
\newcommand{\keywords}[1]{\par\smallskip\noindent
  {\footnotesize\textsc{Keywords} --- \textit{#1}}\par}

% --- The critical apparatus -------------------------------------------------

% The Gallica view number, in the margin. This is the anchor that lets a
% disagreement over a formula be settled by looking at *one* image rather than
% a volume of 750. The site uses it to turn the facsimile as the transcription
% is scrolled. A view is one scanned image: usually one side of a leaf.
\newcommand{\page}[1]{%
  \marginnote{\footnotesize\color{gray!70}\textsf{v.\,#1}}%
  \label{page:#1}\ignorespaces}

% The BnF's pencilled foliation, where the leaf carries it and a pass can read
% it: « 198r ». This is what the literature cites — « ff. 198r–208v » — and
% the one line that turns a citation into a view. Recorded, never inferred:
% a folio a pass cannot read is not written.
\newcommand{\folio}[1]{%
  \marginnote{\footnotesize\color{gray!50}\textsf{f.\,#1}}\ignorespaces}

% The modernised reading's coarser counterpart to \page{}: which views a
% section is drawn from.
\newcommand{\pagerange}[2]{%
  \marginnote{\footnotesize\color{gray!70}\textsf{v.\,#1--#2}}%
  \ignorespaces}

% Illegible: never guessed.
\newcommand{\ill}{\textcolor{red!60!black}{[\,\ldots\,]}}

% A doubtful reading: the word is given, and flagged as uncertain.
\newcommand{\uncertain}[1]{\uline{#1}}

% An editorial addition — a missing letter, an implied word.
\newcommand{\add}[1]{\textcolor{blue!50!black}{[#1]}}

% Struck out by the writer. What was crossed out is part of the document.
\newcommand{\struck}[1]{\sout{#1}}

% The transcriber's note, on the state of the page or the mathematics.
\newcommand{\note}[1]{\par\smallskip{\small\color{gray!60!black}\textit{[#1]}}\par\smallskip}

% A marginal note of hers, distinct from the body of the text.
\newcommand{\marginal}[1]{\par\smallskip\noindent\hspace{1em}%
  {\small\color{gray!60!black}#1}\par\smallskip}

% --- Title ------------------------------------------------------------------
% The title block is French, like the documents themselves.

\AddToHook{shipout/background}{%
  \ifx\grothWatermark\empty\else
    \begin{tikzpicture}[remember picture, overlay]
      \node[rotate=52, scale=3.4, align=center, text=gray!18,
            font=\sffamily\bfseries]
        at (current page.center) {\grothWatermark};
    \end{tikzpicture}%
  \fi}

\AtBeginDocument{%
  \begin{center}
    {\large\bfseries Manuscrits de mathématiciens (Gallica) ---
      \ifx\grothShelfmark\empty\grothFolder\else\grothShelfmark\fi}\\[2pt]
    {\itshape\grothFoldertitle}\\[4pt]
    {\small vues \grothFirst--\grothLast
      \ifx\grothBatch\empty\else\ (lot \grothBatch)\fi}\\[2pt]
    \ifx\grothDating\empty\else
      {\small\color{gray!60!black}Datation du catalogue : \grothDating}\\[2pt]
    \fi
    \ifx\grothWatermark\empty\else
      {\small\bfseries\color{red!45!gray}\grothWatermark}\\[2pt]
    \fi
    {\footnotesize\color{gray!60!black}Transcription établie d'après les images IIIF de Gallica.
     Source gallica.bnf.fr / Bibliothèque nationale de France.\\
     Les lectures incertaines sont soulignées, les illisibles notées [\,\ldots\,],
     les ajouts entre crochets ; \textsf{v.} numérote les vues de Gallica,
     \textsf{f.} la foliotation au crayon de la BnF.}
  \end{center}
  \vspace{1em}}

\endinput


\folder{fr-9115}
\shelfmark{Français 9115}
\batch{12}
\pages{221}{240}
\dating{1801-1900}
\watermark{Lecture automatique\\non vérifiée}
\foldertitle{Papiers de Sophie GERMAIN. — Recueil de dissertations et problèmes mathématiques et physiques.}

\begin{document}

\note{Numérotation des feuillets. Chaque feuillet écrit porte deux nombres à
l'encre : au milieu du haut de la page, souligné, la pagination de la main
elle-même pour cette pièce — 4 (vue 224), 5 (226), 6 (228), 7 (230), 8 (234),
9 (236), 10 (238), 11 (240), le chiffre de la vue 222 étant coupé par le bord
replié du feuillet et lu 3 avec doute, et la vue 232 n'en portant aucun ; en
haut à droite, la série à l'encre des lots précédents, 108 (vue 222), 109
(224), 110 (226), 111 (228), 112 (230), 113 (232), 114 (234), 115 (236), 116
(238), 117 (240), un nombre par feuillet et sans lacune. Tous ces nombres ont
été lus sur la vue ; aucun n'est déduit. Aucune foliotation au crayon de la
Bibliothèque n'a été lue sur ces vues, et aucun « f. » n'est donc porté ici.
Les vues 221, 223, 225, 227, 229, 231, 233, 235, 237 et 239 sont les dos
blancs des feuillets photographiés aux vues précédentes : ce qu'on y croit
lire est le recto vu en transparence, à rebours, et elles ne sont pas
transcrites. La citation, dont les lignes sont toutes ouvertes d'un guillemet
à la manière du temps, commence avant ce lot et se ferme à la vue 228 ; les
guillemets sont conservés en tête des lignes de prose, et non répétés devant
les équations mises hors du texte.}

\page{222}

\note{au milieu du haut de la page, le nombre \uncertain{3}, souligné : le
haut du chiffre passe sous le bord replié du feuillet et n'est pas entier ;
le nombre 108 est à l'encre en haut à droite. Les guillemets ouvrants sont
répétés en tête de chaque ligne de la page, lignes d'équations comprises, à
la manière du temps ; ils sont conservés en tête des lignes de prose, et non
répétés devant les équations mises ici hors du texte. Le verso de ce
feuillet, vue 223, est blanc.}

\[ \text{I.}\ \alpha e^{\lambda\omega}+\beta e^{-\lambda\omega}+\gamma\,\text{sin.}\,\lambda\omega+\delta\cos.\lambda\omega=0 \]
\[ \text{II.}\ \alpha' e^{\lambda\omega}+\beta' e^{-\lambda\omega}+\gamma'\,\text{sin.}\,\lambda\omega+\delta'\cos.\lambda\omega=0 \]
\begin{align*}
&\text{III.}\ \alpha e^{\lambda\omega}-\beta e^{-\lambda\omega}+\gamma\cos.\lambda\omega-\delta\,\text{sin.}\,\lambda\omega=\alpha' e^{\lambda\omega}-\beta' e^{-\lambda\omega}+\gamma'\cos.\lambda\omega-\delta'\,\text{sin.}\,\lambda\omega \\
&\text{ou III}\ (\alpha-\alpha')e^{\lambda\omega}-(\beta-\beta')e^{-\lambda\omega}+(\gamma-\gamma')\cos.\lambda\omega-(\delta-\delta')\,\text{sin.}\,\lambda\omega=0 \\
&\text{IV}\ (\alpha-\alpha')e^{\lambda\omega}+(\beta-\beta')e^{-\lambda\omega}-(\gamma-\gamma')\,\text{sin.}\,\lambda\omega-(\delta-\delta')\cos.\lambda\omega=0
\end{align*}

« Maintenant donc en considérant les extrémités, il en résultera « quatre
équations selon que l'on fera ou $\lambda=0$ ou $\lambda=1$ ; car l'extrémité
« $E$, pour laquelle $u=0$ donne. V. $\alpha+\beta+\delta=0$ et VI.
$\alpha+\beta-\delta=0$ et « l'autre extrémité $E$, pour laquelle $u=1$ donne
VII. $\alpha' e^{\omega}+\beta' e^{-\omega}+\gamma'\,\text{sin.}\,\omega+\delta'\cos.\omega=0$
« et VIII. $\alpha' e^{\omega}+\beta' e^{-\omega}-\gamma'\,\text{sin.}\,\omega-\delta'\cos.\omega=0$.
\note{devant « Maintenant », après le guillemet, un « e » isolé, tracé puis
laissé, sans rature. La seconde extrémité est écrite « $E$ » comme la
première, du même tracé à trois traits ; c'est le $F$ du mémoire. Le mot
« extrémité » est traversé d'un long trait horizontal, qui est la barre du
$t$ rejetée sur le mot et non une rature.}

« Ainsi nous avons rencontré huit équations pour déterminer « les huit
coefficiens $\alpha$, $\beta$, $\gamma$, $\delta$, et $\alpha'$, $\beta'$,
$\gamma'$, $\delta'$ et il faudra « trouver de plus, une équation d'où on
puisse tirer la valeur d'$\omega$.

« Nous commençons par les équations V et VI dont la comparaison « donne
$\beta=-\alpha$ et $\delta=0$ ; ensuite les équations VII et VIII ajoutées
« l'une a l'autre donnent encore $\alpha' e^{\omega}+\beta' e^{-\omega}=0$,
et soustraites « l'une de l'autre elles donnent encore
$\gamma'\,\text{sin.}\,\omega+\delta'\cos.\omega=0$, d'où il « résulte
$\beta'=-\alpha' e^{2\omega}$ et $\delta'=-\gamma'\,\text{tang.}\,\omega$.
Ces valeurs substituées dans « la première et seconde équation donnent

\[ \text{I.}\ \alpha(e^{\lambda\omega}-e^{-\lambda\omega})+\gamma\,\text{sin.}\,\lambda\omega=0 \text{ et II. } \alpha'(e^{\lambda\omega}-e^{\omega(2-\lambda)})+\gamma'\,\text{sin.}\,\lambda\omega-\gamma'\cos.\lambda\omega\,\text{tang.}\,\omega=0. \]

« De ces équations on tire les valeurs :
$\gamma=-\alpha(e^{\lambda\omega}-e^{-\lambda\omega})$ et
$\gamma'=-\dfrac{\alpha'(e^{\lambda\omega}-e^{\omega(2-\lambda)})}{\text{sin.}\,\lambda\omega-\cos.\lambda\omega\,\text{tang.}\,\omega}$
« et par conséquent
$\delta'=\dfrac{\alpha'(e^{\lambda\omega}-e^{\omega(2-\lambda)})\,\text{tang.}\,\omega}{\text{sin.}\,\lambda\omega-\cos.\lambda\omega\,\text{tang.}\,\omega}$.
\note{la valeur de $\gamma$ est écrite sans dénominateur ; la ligne qui donne
plus bas $\gamma-\gamma'$ porte bien $\text{sin.}\,\lambda\omega$ sous le
même numérateur.}

« en rassemblant ces valeurs pour les autres équations nous avons.

\begin{align*}
&\beta-\beta'=-\alpha+\alpha' e^{2\omega},\quad \gamma-\gamma'=-\frac{\alpha(e^{\lambda\omega}-e^{-\lambda\omega})}{\text{sin.}\,\lambda\omega}+\frac{\alpha'(e^{\lambda\omega}-e^{\omega(2-\lambda)})}{\text{sin.}\,\lambda\omega-\cos.\lambda\omega\,\text{tang.}\,\omega} \\
&\delta-\delta'=-\frac{\alpha'(e^{\lambda\omega}-e^{\omega(2-\lambda)})\,\text{tang.}\,\omega}{\text{sin.}\,\lambda\omega-\cos.\lambda\omega\,\text{tang.}\,\omega}.
\end{align*}

« on a maintenant

\begin{align*}
&\text{III.}\ \alpha\big\{(e^{\lambda\omega}+e^{-\lambda\omega})-(e^{\lambda\omega}-e^{-\lambda\omega})\cot.\lambda\omega\big\} \\
&\qquad=\alpha'\Big\{(e^{\lambda\omega}+e^{\omega(2-\lambda)})-(e^{\lambda\omega}-e^{\omega(2-\lambda)})\frac{\cos.\lambda\omega+\text{sin.}\,\lambda\omega\,\text{tang.}\,\omega}{\text{sin.}\,\lambda\omega-\cos.\lambda\omega\,\text{tang.}\,\omega}\Big\}
\end{align*}

« équation qui peut être mise sous cette forme :

\begin{align*}
&\alpha\big\{(e^{\lambda\omega}+e^{-\lambda\omega})-(e^{\lambda\omega}-e^{-\lambda\omega})\cot.\lambda\omega\big\} \\
&\qquad=\alpha'\big\{(e^{\lambda\omega}+e^{\omega(2-\lambda)})-(e^{\lambda\omega}-e^{\omega(2-\lambda)})\big\}\cot(1-\lambda)\omega
\end{align*}

« l'équation quatrième donne celle-ci
$2\alpha(e^{\lambda\omega}-e^{-\lambda\omega})=2\alpha'(e^{\lambda\omega}-e^{\omega(2-\lambda)})$ ;
\note{dans la forme réduite, le second membre porte un « $-$ » devant
$(e^{\lambda\omega}-e^{\omega(2-\lambda)})$, et l'accolade se ferme avant
$\cot(1-\lambda)\omega$, qui reste hors de l'accolade : c'est ainsi que la
ligne est écrite. Le feuillet s'arrête sur cette équation, au tiers
supérieur de la page ; le reste est blanc.}

\page{224}

\note{au milieu du haut de la page, le nombre 4, souligné ; le nombre 109 est
à l'encre en haut à droite. Comme à la vue précédente, le guillemet ouvrant
est répété en tête de chaque ligne, titre compris. Le verso, vue 225, est
blanc.}

« et on déduit de cette dernière
$\dfrac{\alpha}{\alpha'}=\dfrac{e^{\lambda\omega}-e^{\omega(2-\lambda)}}{e^{\lambda\omega}-e^{-\lambda\omega}}$
« d'où on conclut qu'il est permis de prendre
$\alpha=e^{\lambda\omega}-e^{\omega(2-\lambda)}$, et
$\alpha'=e^{\lambda\omega}-e^{-\lambda\omega}$ « valeurs qu'il faut
substituer maintenant dans la troisième équation ; « cette substitution
donne après les réductions convenables

\begin{align*}
&0=2-2e^{2\omega}-(e^{\lambda\omega}-e^{\omega(2-\lambda)})(e^{\lambda\omega}-e^{-\lambda\omega})(\cot.\lambda\omega+\cot.(1-\lambda)\omega) \quad\text{ou} \\
&0=2-2e^{2\omega}-(e^{\lambda\omega}-e^{\omega(2-\lambda)})(e^{\lambda\omega}-e^{-\lambda\omega})\frac{\text{sin.}\,\omega}{\text{sin.}\,\lambda\omega\,\text{sin.}(1-\lambda)\omega}
\end{align*}

« ainsi la question est réduite de cette équation les valeurs de l'angle
« $\omega$, je n'ose entreprendre ce travail a cause de la complication des
« formules, il suffit d'avoir montré dans ce lieu par quelle méthode « on
peut traiter ces difficiles questions \ldots\ mais le « cas ou
$\lambda=\frac{1}{2}$ mérite d'être développé.
\note{« la question est réduite de cette équation les valeurs de l'angle
$\omega$ » : la ligne est écrite ainsi, sans le verbe qui vaudrait « à tirer
de ». Après « questions », neuf points de suite tiennent la place de ce qui
n'est pas traduit.}

\section{« Développement du cas « dans lequel la lame $E$ $F$ est appuiée non seulement a ses extrémités $E$ et $F$, mais « encor dans le milieu de dela longueur}

\note{le titre tient trois lignes, chacune ouverte d'un guillemet ; un trait
horizontal le sépare de la suite. « de dela longueur » est écrit ainsi, avec
« de » puis son « dela » en un mot.}

« §. 48. Parceque dans ce cas $\lambda=\frac{1}{2}$, l'équation finale prend
cette forme :
\[ 0=2-2e^{2\omega}-(e^{\frac{1}{2}\omega}-e^{\frac{3}{2}\omega})(e^{\frac{1}{2}\omega}-e^{-\frac{1}{2}\omega})\frac{\text{sin.}\,\omega}{\text{sin.}^2\frac{1}{2}\omega} \]
« qui peut être réduite a celle ci : \struck{\ill{}}
\[ 0=2(1+e^{\omega})\,\text{sin.}^2\tfrac{1}{2}\omega-(e^{\omega}-1)\,\text{sin.}\,\omega, \]
« en écrivant $2\,\text{sin.}\,\frac{1}{2}\omega\cos.\frac{1}{2}\omega$ au
lieu de $\text{sin.}\,\omega$ « cette équation se change en celle-ci
\[ 0=(1+e^{\omega})\,\text{sin.}^2\tfrac{1}{2}\omega-(e^{\omega}-1)\,\text{sin.}\,\tfrac{1}{2}\omega\cos.\tfrac{1}{2}\omega \]
\note{entre « réduite a celle ci : » et l'équation qui suit, une formule est
entièrement biffée d'un zigzag épais ; on n'y distingue plus que des
exposants $2\omega$, $\frac{1}{2}\omega$ et $\frac{3}{2}\omega$, et elle
n'est pas lue. Le facteur écrit ici est $2(1+e^{\omega})$, l'exposant est
bien $\omega$.}

« laquelle a manifestement deux facteurs l'un
$\text{sin.}\,\frac{1}{2}\omega$ et l'autre «
$(1+e^{\omega})\,\text{sin.}\,\frac{1}{2}\omega-(e^{\omega}-1)\cos.\frac{1}{2}\omega$,
qui égalés a zéro donnent chacun une « solution : nous les examinerons donc
chacun en particulier.

« §. \uncertain{49}. Nous prenons pour le premier cas
$\text{sin.}\,\frac{1}{2}\omega=0$, et on aura par conséquent «
$\frac{1}{2}\omega=i\pi$, $i$ représentant un nombre entier quelconque, ce
qui comprend déjà « une quantité innombrable de mouvemens réguliers ; il est
clair que « ce cas est parfaitement semblable au cas quatrième que nous «
avons developpé plus haut [ce cas est celui ou les deux extrémités « sont
appuiées], si ce n'est qu'on a ici $\frac{1}{2}\omega$ a la place d'$\omega$
« dans l'endroit cité, et cela parceque l'une et l'autre parties « dela
longueur sont seulement la moitié de $a$.
\note{le chiffre du paragraphe est brouillé par une reprise à la plume ; il
est lu 49 avec doute. Les crochets droits sont de la main du feuillet. Le
feuillet s'arrête ici, un peu au-dessous du milieu de la page.}

\page{226}

\note{au milieu du haut de la page, le nombre 5, souligné ; le nombre 110 est
à l'encre en haut à droite. Guillemet ouvrant en tête de chaque ligne. Le
verso, vue 227, est blanc.}

« L'une et l'autre moitiés font leurs vibrations, de la même manière « que si
elles avoient leurs extrémités appuiées au point ou est le stilet « c'est
pourquoi les sons simples que l'une et l'autre portions « peuvent donner sont
exprimés par les nombres suivans:

\[ \frac{4\pi c\surd 2gb}{aa};\ \frac{16\pi c\surd 2gb}{aa};\ \frac{36\pi c\surd 2gb}{aa};\ \frac{64\pi c\surd 2gb}{aa}; \]

« et sont par conséquent les doubles octaves audessus de ceux qui ont « lieu
dans le cas quatrième, la raison en est que dans ce cas « les longueurs sont
seulement lamoitié de celles de l'autre cas

« §. 50. Ainsi dans ce cas les coëfficiens $\alpha$, $\beta$, $\gamma$,
$\delta$ et $\alpha'$, $\beta'$, $\gamma'$, $\delta'$ seront « déterminés
dela manière suivante :

\begin{align*}
&\alpha=e^{\frac{1}{2}\omega}-e^{\frac{3}{2}\omega};\quad \beta=e^{\frac{3}{2}\omega}-e^{\frac{1}{2}\omega};\quad \gamma=-\frac{(e^{\frac{1}{2}\omega}-e^{\frac{3}{2}\omega})(e^{\frac{1}{2}\omega}-e^{-\frac{1}{2}\omega})}{\text{sin.}\,\frac{1}{2}\omega};\quad \delta=0 \\
&\alpha'=e^{\frac{1}{2}\omega}-e^{-\frac{1}{2}\omega};\quad \beta'=-e^{2\omega}(e^{\frac{1}{2}\omega}-e^{-\frac{1}{2}\omega});\quad \gamma'=-\frac{(e^{\frac{1}{2}\omega}-e^{-\frac{1}{2}\omega})(e^{\frac{1}{2}\omega}-e^{\frac{3}{2}\omega})}{\text{sin.}\,\frac{1}{2}\omega-\cos.\frac{1}{2}\omega\,\text{tang.}\,\omega}; \\
&\delta'=\frac{(e^{\frac{1}{2}\omega}-e^{-\frac{1}{2}\omega})(e^{\frac{1}{2}\omega}-e^{\frac{3}{2}\omega})\,\text{tang.}\,\omega}{\text{sin.}\,\frac{1}{2}\omega-\cos.\frac{1}{2}\omega\,\text{tang.}\,\omega}
\end{align*}

« or,
$\text{sin.}\,\frac{1}{2}\omega-\cos.\frac{1}{2}\omega\,\text{tang.}\,\omega=\dfrac{\text{sin.}\,\frac{1}{2}\omega\cos.\omega-\cos.\frac{1}{2}\omega\,\text{sin.}\,\omega}{\cos.\omega}=-\dfrac{\text{sin.}\,\frac{1}{2}\omega}{\cos.\omega}$,
a cause de « $\omega=2i\pi$, on a $\cos.\omega=1$ et
$\text{sin.}\,\frac{1}{2}\omega=0$. Si nous multiplions donc tous ces «
coëfficiens par $\text{sin.}\,\frac{1}{2}\omega$ nous aurons $\alpha=0$,
$\beta=0$, $\gamma=-(1-e^{\omega})(e^{\omega}-1)$, $\delta=0$ ; « enfin
$\alpha'=0$, $\beta'=0$, $\gamma'=(e^{\omega}-1)(1-e^{\omega})$ et
$\delta'=0$.

« Ainsi parceque tous ces coëfficiens s'évanouissent, excepté $\gamma$ et
$\gamma'$ « et que de plus $\gamma=-\gamma'$ si nous faisons $\gamma=1$ nous
aurons $\gamma'=-1$ ; « ainsi pour la portion $E L$ l'équation qui exprime le
mouvement « sera
\[ y=C\,\text{sin.}\Big(\zeta+t.\frac{\omega\omega c\surd 2gb}{aa}\Big)\text{sin.}\,u\omega \]
et pour la portion $L F$
\[ y=-C\,\text{sin.}\Big(\zeta+t.\frac{\omega\omega c\surd 2gb}{aa}\Big)\text{sin.}\,u\omega \]
dans lesquelles $s=\frac{u}{a}$,
\note{devant la seconde équation, trois petits traits en ligne comme des
points de conduite, à la place du facteur laissé de côté. La lettre posée
devant le signe d'égalité, à « dans lesquelles », a la forme d'un long
« s » ; le rapport lui-même est écrit $\frac{u}{a}$. Laissé tel que lu.}

« §. 51. Mais outre cette solution on en tire une autre de l'autre facteur «
duquel il résulte
\[ \text{tang.}\,\tfrac{1}{2}\omega=\frac{e^{\omega}-1}{1+e^{\omega}}=\frac{e^{\frac{1}{2}\omega}-e^{-\frac{1}{2}\omega}}{e^{\frac{1}{2}\omega}+e^{-\frac{1}{2}\omega}} \]
« valeur qui est lamême que celle que nous avons trouvé cidessus cas V. «
[ce cas est celui ou une extrémité est fixée et l'autre appuiée] « la seule
différence consiste en ce qu'on a ici $\frac{1}{2}\omega$ au lieu de
$\omega$, ce qui est naturel « puisque l'une et l'autre portions dela
longueur sont seulement lamoitié « de $a$.

« Nous concevons presentement que l'une et l'autre portions $E L$ « et $L F$
puisse se mouvoir en même tems et que les deux extrémités « $E$ et $F$ étant
simplement appuiées, le point $L$ soit fixe. Les « valeurs d'$\omega$ seront
alors $\frac{5\pi}{2}$, $\frac{9\pi}{2}$, $\frac{13\pi}{2}$, \&c. et par
conséquent le son sera « de deux octaves plus élevé que dans le cas V. Il est
digne de remarque « que les deux portions $E L$ et $L F$ peuvent vibrer de
deux manières,
\note{les mots « outre » et « cette », au §. 51, sont traversés d'un long
trait horizontal qui est la barre du $t$ rejetée, non une rature. Les
crochets droits sont de la main du feuillet. « élevé » est tracé comme
« ébré », sans hampe à la deuxième lettre. Le feuillet s'arrête sur
« manières, », un peu au-dessous du milieu de la page.}

\page{228}

\note{au milieu du haut de la page, le nombre 6, souligné ; le nombre 111 est
à l'encre en haut à droite. La citation s'achève à la troisième ligne, sur un
guillemet fermant ; à partir de là le feuillet porte le texte de la main
elle-même, sans guillemets. Le verso, vue 229, est blanc.}

« L'une qui se rapporte au cas IV et l'autre au cas V. Au reste les « valeurs
des coëfficiens sont les mêmes que celles que nous avons trouvées « plus
haut, si ce n'est que pour ce cas on n'a pas
$\text{sin.}\,\frac{1}{2}\omega=0$ ».

I. La première réflexion qui se présente ici est que cette seconde solution
sort de l'état de la question car l'énoncé du problème indique assez que les
mouvemens dela lame soumise a l'action du stilet en $L$ doivent être compris
au nombre de tous les mouvemens possibles dont la même lame est susceptible
lorsqu'elle a des extrémités appuiées et les formules employées a trouver
l'équation deffinitive alaquelle il s'agit de satisfaire confirment cette
remarque ; puisque l'auteur \struck{y employe} a recours a \struck{y} quatre
équations dépendantes de l'état des extrémités.
\note{la page s'ouvre sur un « I. » isolé devant la première ligne, de la
même encre, comme les lettres de série des autres brouillons. « a recours a »
est écrit au-dessus de la ligne, par-dessus les mots biffés.}

Cela posé, il en résulte, ce me semble, que la condition
$\text{sin.}\,\omega=0$, qu'il a établi \struck{\ill{}} devoir toujours
\struck{avoir lieu} être remplie pour tous les mouvemens possibles dela lame
dont il s'agit, doit s'accorder avec la solution du present problème ; et
c'est effectivement ce que l'on trouve pour la première solution savoir
$\text{sin.}\,\frac{1}{2}\omega=0$, mais la même chose n'a pas lieu pour la
seconde solution qu'il admet, car l'équation
$\text{tang.}\,\frac{1}{2}\omega=\frac{e^{\omega}-1}{e^{\omega}+1}$ ne peut
s'accorder avec la condition $\text{sin.}\,\omega=0$ que dans la supposition
$\omega=0$
\marginal{(\ill{} cas IV §. 35 \uncertain{\&c.})}
\note{« être remplie » est écrit au-dessus de « avoir lieu » biffé ; le mot
biffé après « établi » est noyé sous la rature et n'est pas lu. Sous la même
ligne, à droite, un court renvoi entre parenthèses : on y lit « cas IV §.
35 », précédé d'une lettre qui n'est pas décidée et suivi d'un signe lu
« \&c. » avec doute.}

Cependant le respect dû \struck{au \ill{}} l'opinion d'Euler m'a engagé a
examiner le cas general \marginal{$\times$ afin de savoir si il étoit
susceptible d'une solution dans laquelle celle qui m'avoit inspiré des doutes
peut être comprise} En cherchant donc a résoudre son équation
\[ 0=2(1-e^{2\omega})-(e^{\lambda\omega}-e^{\omega(2-\lambda)})(e^{\lambda\omega}-e^{-\lambda\omega})\frac{\text{sin.}\,\omega}{\text{sin.}\,\lambda\omega\,\text{sin.}(1-\lambda)\omega} \]
pour une valeur quelconque de $\lambda$, j'ai voulu m'assurer de l'existence
du facteur
$\text{tang.}\,\lambda\omega=\frac{e^{2\lambda\omega}-1}{e^{2\lambda\omega}+1}$,
pour les valeurs de $\lambda$ differentes $\frac{1}{2}$.
\note{« l'opinion » est écrit au-dessus des mots biffés qui suivent « dû » ;
sous la rature on distingue « au » et un mot qui n'est pas lu. « donc » est
ajouté au-dessus de la ligne, après « cherchant ». La croix qui suit « le cas
general » appelle la note portée dans la marge de gauche, en trois lignes,
qui commence elle-même par la même croix ; elle est mise ici à sa place.
Le feuillet s'arrête sur « $\frac{1}{2}$. » ; le reste de la page est
blanc.}

\page{230}

\note{au milieu du haut de la page, le nombre 7, souligné ; le nombre 112 est
à l'encre en haut à droite. « Pour cela » est écrit seul au-dessus de la
première ligne, un peu à gauche. Le verso, vue 231, est blanc.}

Pour cela

2. Je reprend \struck{donc} l'équation
\[ 0=2(1-e^{2\omega})-(e^{\lambda\omega}-e^{(2-\lambda)\omega})(e^{\lambda\omega}-e^{-\lambda\omega})\frac{\text{sin.}\,\omega}{\text{sin.}\,\lambda\omega\,\text{sin.}(1-\lambda)\omega} \]
ou
\[ 0=2(1-e^{2\omega})\,\text{sin.}\,\lambda\omega\,\text{sin.}(1-\lambda)\omega-(1-e^{2(1-\lambda)\omega})(e^{2\lambda\omega}-1)\,\text{sin.}\,\omega \]
\note{à la fin de la première de ces deux équations, un « $=0$ » pâle, à demi
effacé, après la fraction ; la ligne commence déjà par « $0=$ ». Le mot
« donc », après « reprend », est traversé d'un trait épais et court.}

Si on veut conserver la plus parfaite analogie entre le cas général et celui
ou $\lambda=\frac{1}{2}$, on peut diviser cette équation par le facteur
$e^{2\lambda\omega}-1$, duquel on tireroit, comme dans le cas cité
$\omega=0$. L'équation elle se réduit par cette division a
\begin{align*}
&2\big\{1+e^{2\lambda\omega}+e^{4\lambda\omega}+e^{6\lambda\omega}\ldots+e^{2(1-\lambda)\omega}\big\}\,\text{sin.}\,\lambda\omega\,\text{sin.}(1-\lambda)\omega \\
&\qquad -(1-e^{2(1-\lambda)\omega})\,\text{sin.}\,\omega=0
\end{align*}

J'observe que par la nature de laquestion, $\lambda$ est toujours egal aune
fraction et que dans le cas des deux extrémités appuiées, cas dont il s'agit
ici, cette fraction est partie aliquote de l'unité ou multiple d'une telle
\struck{fraction} partie.
\note{« multiple d'une telle fraction partie. » est ajouté au-dessous de la
ligne, en petits caractères, et « fraction » y est biffé.}

On trouve par les formules connues, pour le cas ou le dénominateur de la
fraction $\lambda$ est un nombre impair et supposant le numerateur egal a
l'unité,
\[ \text{sin.}\,\omega=\frac{1}{\lambda}\,\text{sin.}\,\lambda\omega-\frac{1-\lambda^2}{\lambda^3.2.3}\,\text{sin.}^3\lambda\omega+\frac{(1-\lambda^2)(1-9\lambda^2)}{\lambda^5.\,2.3.4.5}\,\text{sin.}^5\lambda\omega-\text{\&c.} \]
et pour celui ou le dénominateur dela même fraction $\lambda$ est un nombre
pair
\[ \text{sin.}\,\omega=\cos.\lambda\omega\Big\{\frac{1}{\lambda}\,\text{sin.}\,\lambda\omega-\frac{1-4\lambda^2}{\lambda^3.2.3}\,\text{sin.}^3\lambda\omega+\frac{(1-4\lambda^2)(1-16\lambda^2)}{\lambda^5.\,2.3.4.5}\,\text{sin.}^5\lambda\omega\ \text{\&c}\Big\} \]
\note{« ou » est écrit au-dessus de « et supposant » ; un long trait
horizontal court sur « et supposant le numerateur », qui peut être la barre
du $t$ rejetée et n'est pas donné pour une rature.}

L'équation qu'il s'agit de resoudre se réduit donc a l'une ou l'autre des
formes suivantes
\begin{align*}
&2\big\{1+e^{2\lambda\omega}+e^{4\lambda\omega}+e^{6\lambda\omega}\ldots+e^{2(1-\lambda)\omega}\big\}\,\text{sin.}\,\lambda\omega\,\text{sin.}(1-\lambda)\omega \\
&\quad -(1-e^{2(1-\lambda)\omega})\Big\{\frac{1}{\lambda}-\frac{1-\lambda^2}{\lambda^3.2.3}\,\text{sin.}^2\lambda\omega+\frac{(1-\lambda^2)(1-9\lambda^2)}{\lambda^5.\,2.3.4.5}\,\text{sin.}^4\lambda\omega-\text{\&c.}\Big\}\,\text{sin.}\,\lambda\omega=0
\end{align*}
\begin{align*}
&2\big\{1+e^{2\lambda\omega}+e^{4\lambda\omega}+e^{6\lambda\omega}+e^{2(1-\lambda)\omega}\big\}\,\text{sin.}\,\lambda\omega\,\text{sin.}(1-\lambda)\omega \\
&\quad -(1-e^{2(1-\lambda)\omega})\cos.\lambda\omega\Big\{\frac{1}{\lambda}-\frac{1-4\lambda^2}{\lambda^3.2.3}\,\text{sin.}^2\lambda\omega+\frac{(1-4\lambda^2)(1-16\lambda^2)}{\lambda^5.\,2.3.4.5}\,\text{sin.}^4\lambda\omega\Big\}\,\text{sin.}\,\lambda\omega=0
\end{align*}

ainsi dans l'un et l'autre cas $\text{sin.}\,\lambda\omega$ est un des
facteurs de cette équation et en egalant ce facteur a zéro on a une solution
\struck{\ill{}} qui se rapporte au cas IV et delaquelle il résulte
$\omega=\frac{i\pi}{\lambda}$, $i$ étant, comme dans le cas
$\lambda=\frac{1}{2}$, \struck{\uncertain{examiné} \ill{}}, un nombre entier
quelconque. on en conclut \struck{de même} comme dans ce dernier cas, en
observant que la lame se partage alors en $\frac{1}{\lambda}$ parties de
lalongueur \struck{$\lambda a$}, que ces differentes parties font leurs
vibrations dela même manières que si elles étoient autant delames
differentes et de pareille longueur ayant leurs extremités appuiées tant dans
le point ou est placé le stilet que dans les autres points de repos, que l'on
pourroit nommer sympatiques,
\note{« ou » est ajouté au-dessus de « le cas $\lambda=\frac{1}{2}$ » ; les
deux ou trois mots biffés qui suivent ne se laissent pas lire sous la rature,
le premier ressemble à « examiné ». « comme dans ce dernier cas » est écrit
au-dessus de « de même » biffé. Le feuillet s'arrête sur « sympatiques, » ;
le bas de la page est blanc.}

\page{232}

\note{feuillet plus grand que les précédents, monté sur onglet ; le nombre
113 est à l'encre en haut à droite, et cette page ne porte pas de nombre au
milieu du haut, à la différence des vues 222 à 230 et 234 à 240. Le texte
n'occupe que le cinquième supérieur de la page ; tout le reste est blanc. Sa
première ligne ne fait pas suite à la dernière de la vue 230. Le verso, vue
233, est blanc.}

est satisfaite ici puisqu'elle résulte dela supposition
$\text{sin.}\,\lambda\omega=0$.

2 Cela posé, on peut encore suivre, dans la détermination des coëfficiens
$\alpha$, $\beta$, $\gamma$, $\delta$ et $\alpha'$, $\beta'$, $\gamma'$,
$\delta'$, l'analogie entre le cas général et celui qu'Euler a calculé car on
trouve dela même manière
\begin{align*}
&\alpha=e^{\lambda\omega}-e^{(2-\lambda)\omega};\quad \beta=e^{(2-\lambda)\omega}-e^{\lambda\omega};\quad \gamma=\frac{(e^{(2-\lambda)\omega}-e^{\lambda\omega})(e^{\lambda\omega}-e^{-\lambda\omega})}{\text{sin.}\,\lambda\omega};\quad \delta=0 \\
&\alpha'=e^{\lambda\omega}-e^{-\lambda\omega};\quad \beta'=-e^{2\omega}(e^{\lambda\omega}-e^{-\lambda\omega});\quad \gamma'=\frac{(e^{(2-\lambda)\omega}-e^{\lambda\omega})(e^{\lambda\omega}-e^{-\lambda\omega})}{\text{sin.}\,\lambda\omega-\cos.\lambda\omega\,\text{tang.}\,\omega}; \\
&\delta'=\frac{(e^{\lambda\omega}-e^{(2-\lambda)\omega})(e^{\lambda\omega}-e^{-\lambda\omega})\,\text{tang.}\,\omega}{\text{sin.}\,\lambda\omega-\cos.\lambda\omega\,\text{tang.}\,\omega}
\end{align*}

et a cause de $\text{tang.}\,\omega=0$ on peut écrire
\[ \gamma'=\frac{(e^{(2-\lambda)\omega}-e^{\lambda\omega})(e^{\lambda\omega}-e^{-\lambda\omega})}{\text{sin.}\,\lambda\omega};\qquad \delta'=\frac{0}{\text{sin.}\,\lambda\omega} \]
si on multiplie donc tous ces coëfficiens par $\text{sin.}\,\lambda\omega$ on
aura comme dans le cas ou $\lambda=\frac{1}{2}$

\struck{$\alpha=0$, $\beta=0$, $\gamma=$}

a cause de
$\text{sin.}\,\lambda\omega-\cos.\lambda\omega\,\text{tang.}\,\omega=-\dfrac{\text{sin.}(1-\lambda)\omega}{\cos.\omega}$
on peut écrire en observant que d'où il résulte
\note{la ligne courte « $\alpha=0$, $\beta=0$, $\gamma=$ » est biffée d'un
trait épais. Dans la ligne qui commence par « a cause de », un trait
horizontal fin court sur toute la longueur ; il n'est pas donné pour une
rature, la ligne suivante reprenant la même phrase. « d'où il résulte » est
écrit au-dessous de cette ligne, en petits caractères, entre « de » et
« $\text{sin.}\,\lambda\omega$ ».}

$+$ a cause de $\cos.\omega=\uncertain{1}$ et de
$\text{sin.}\,\lambda\omega-\cos.\lambda\omega\,\text{tang.}\,\omega=-\text{sin.}(1-\lambda)\omega$
on peut écrire
\[ \gamma'=-\frac{(e^{(2-\lambda)\omega}-e^{\lambda\omega})(e^{\lambda\omega}-e^{-\lambda\omega})}{\text{sin.}(1-\lambda)\omega};\qquad \delta'=\frac{(e^{\lambda\omega}-e^{(2-\lambda)\omega})(e^{\lambda\omega}-e^{-\lambda\omega})}{\text{sin.}(1-\lambda)\omega} \]
si on multiplie donc les coëfficiens $\alpha$, $\beta$, $\gamma$, $\delta$
par $\text{sin.}\,\lambda\omega$ et les coëfficiens $\alpha'$, $\beta'$,
$\gamma'$, $\delta'$ par $\text{sin.}(1-\lambda)\omega$ ce qui revient au
même puisque ces deux sinus sont l'un et l'autre egaux a zero on aura comme
dans le cas ou $\lambda=\frac{1}{2}$
\note{la ligne s'ouvre sur une croix ou un « + ». Le chiffre qui suit
« $\cos.\omega=$ » est une tache d'encre pleine, d'un seul trait épaissi ; le
calcul qui suit demande 1, et c'est la lecture offerte, avec doute. « donc »
est ajouté au-dessus de « si on multiplie ». Le feuillet s'arrête sur
« $\lambda=\frac{1}{2}$ » ; les quatre cinquièmes inférieurs de la page sont
blancs.}

\page{234}

\note{au milieu du haut de la page, le nombre 8, souligné ; le nombre 114 est
à l'encre en haut à droite. La page reprend, en les développant, les deux
premiers alinéas du feuillet de la vue 232, qui en est une première mise au
net restée sans numéro de page. Le verso, vue 235, est blanc.}

et par conséquent que les sons simples que les parties fractionnaires dela
lame peuvent donner sont exprimés par les nombres suivans:
\[ \frac{\pi c\surd 2gb}{\lambda\lambda aa};\ \frac{4\pi c\surd 2gb}{\lambda\lambda aa};\ \frac{9\pi c\surd 2gb}{\lambda\lambda aa};\ \frac{16\pi c\surd 2gb}{\lambda\lambda aa};\ \text{\&c.} \]

Il est clair que la condition $\text{sin.}\,\omega=0$, qui par la nature dela
question doit avoir lieu en même tems que celle qui résulte dela solution de
l'équation du problème, est satisfaite ici puisqu'elle résulte de la
supposition $\text{sin.}\,\lambda\omega$.
\note{la phrase s'arrête sur « $\text{sin.}\,\lambda\omega$ » sans le
« $=0$ » que porte la même phrase à la vue 232.}

2 Cela posé, on peut suivre encore l'analogie entre le cas général et celui
qu'Euler a calculé dans la détermination des coëfficiens $\alpha$, $\beta$,
$\gamma$ $\delta$ et $\alpha'$, $\beta'$, $\gamma'$, $\delta'$ car on trouve
dela même manière
\begin{align*}
&\alpha=e^{\lambda\omega}-e^{(2-\lambda)\omega};\quad \beta=e^{(2-\lambda)\omega}-e^{\lambda\omega};\quad \gamma=\frac{(e^{(2-\lambda)\omega}-e^{\lambda\omega})(e^{\lambda\omega}-e^{-\lambda\omega})}{\text{sin.}\,\lambda\omega};\quad \delta=0 \\
&\alpha'=e^{\lambda\omega}-e^{-\lambda\omega};\quad \beta'=-e^{2\omega}(e^{\lambda\omega}-e^{-\lambda\omega});\quad \gamma'=\frac{(e^{(2-\lambda)\omega}-e^{\lambda\omega})(e^{\lambda\omega}-e^{-\lambda\omega})}{\text{sin.}\,\lambda\omega-\cos.\lambda\omega\,\text{tang.}\,\omega}; \\
&\delta'=\frac{(e^{\lambda\omega}-e^{(2-\lambda)\omega})(e^{\lambda\omega}-e^{-\lambda\omega})\,\text{tang.}\,\omega}{\text{sin.}\,\lambda\omega-\cos.\lambda\omega\,\text{tang.}\,\omega}.
\end{align*}
et on peut écrire, a cause de $\text{tang.}\,\omega=0$,
\[ \gamma'=\frac{(e^{(2-\lambda)\omega}-e^{\lambda\omega})(e^{\lambda\omega}-e^{-\lambda\omega})}{\text{sin.}\,\lambda\omega},\qquad \text{et}\ \delta'=\frac{(e^{\lambda\omega}-e^{(2-\lambda)\omega})(e^{\lambda\omega}-e^{-\lambda\omega})}{\text{sin.}\,\lambda\omega} \]
si on multiplie donc tous ces coëfficiens par $\text{sin.}\,\lambda\omega$,
on aura comme dans le cas ou $\lambda=\frac{1}{2}$
\begin{align*}
&\alpha=0,\ \beta=0,\ \gamma=-(e^{\lambda\omega}-e^{(2-\lambda)\omega})(e^{\lambda\omega}-e^{-\lambda\omega}),\ \delta=0 \\
&\alpha'=0,\ \beta'=0,\ \gamma'=(e^{\lambda\omega}-e^{(2-\lambda)\omega})(e^{\lambda\omega}-e^{-\lambda\omega}),\ \delta'=0
\end{align*}
et par conséquent $\gamma=-\gamma'$, on peut faire également ici
\[ \gamma=1\qquad \gamma'=-1 \]

Ainsi \marginal{$+$ en prenant pour exemple le cas representé par la fig. 7}
et supposant le stilet placé en $L$ l'équation qui détermine le mouvement
dela partie $E L$ sera
\[ y=C\,\text{sin.}\Big(\zeta+t.\frac{\omega\omega c\surd 2gb}{aa}\Big)\text{sin.}\,\lambda\omega, \]
et pour la partie $L E'$ en tout semblable a la partie $L L'$
on aura
\[ y=-C\,\text{sin.}\Big(\zeta+t.\frac{\omega\omega c\surd 2gb}{aa}\Big)\text{sin.}\,\lambda\omega \]
il est clair que le mouvement dela partie $L' L''$ est déterminé de même par
l'équation
\[ y=C\,\text{sin.}\Big(\zeta+t.\frac{\omega\omega c\surd 2gb}{aa}\Big)\text{sin.}\,\lambda\omega, \]
celui de la partie $L'' L'''$ par l'équation
\[ y=-C\,\text{sin.}\Big(\zeta+t.\frac{\omega\omega c\surd 2gb}{aa}\Big)\text{sin.}\,\lambda\omega \]
et ainsi de suite en adoptant le signe $+$ pour les parties dela courbe qui
sont situées audessus de l'axe et le signe $-$ pour celles qui sont placées
audessous du même axe ou ligne de repos.
\note{la croix qui suit « Ainsi » appelle la note portée dans la marge de
gauche, en deux lignes, qui s'ouvre elle-même sur la même croix ; elle est
mise ici à sa place. Dans « la partie $L E'$ », la seconde lettre est tracée comme le $E$
de « $E L$ » de la ligne précédente, avec un accent, et deux petits traits
sont ajoutés au-dessus des deux lettres ; le reste de la page nomme les
parties $L L'$, $L' L''$, $L'' L'''$, mais la vue 240 écrit bien « $L E'$ »
pour la portion qui va du stilet à l'autre bout. Laissé tel que lu. Le
feuillet s'arrête sur « ligne de repos. » ; le bas de la page est blanc.}

\page{236}

\note{au milieu du haut de la page, le nombre 9, souligné ; le nombre 115 est
à l'encre en haut à droite. Le verso, vue 237, est blanc.}

4. Je reprend apresent l'équation du problème, elle a pour second facteur
\begin{align*}
&2\big\{1+e^{2\lambda\omega}+e^{4\lambda\omega}+e^{6\lambda\omega}\ldots+e^{2(1-\lambda)\omega}\big\}\,\text{sin.}(1-\lambda)\omega \\
&\quad -(1-e^{2(1-\lambda)\omega})\Big\{\frac{\uncertain{2}}{\lambda}-\frac{1-\lambda^2}{2.3.\lambda^3}\,\text{sin.}^2\lambda\omega+\frac{(1-\lambda^2)(1-9\lambda^2)}{2.3.4.5.\,\lambda^5}\,\text{sin.}^4\lambda\omega-\text{\&c.}\Big\}
\end{align*}
ou
\begin{align*}
&2\big\{1+e^{2\lambda\omega}+e^{4\lambda\omega}+e^{6\lambda\omega}\ldots+e^{2(1-\lambda)\omega}\big\}\,\text{sin.}(1-\lambda)\omega \\
&\quad -(1-e^{2(1-\lambda)\omega})\Big\{\frac{1}{\lambda}-\frac{1-4\lambda^2}{2.3.\lambda^3}\,\text{sin.}^2\lambda\omega+\frac{(1-4\lambda^2)(1-16\lambda^2)}{2.3.4.5.\,\lambda^5}\,\text{sin.}^4\lambda\omega\Big\}\cos.\lambda\omega
\end{align*}
suivant que le dénominateur de la fraction $\lambda$ est impair ou pair.
\note{le numérateur de la première fraction de la première accolade est un
petit chiffre à boucle, lu 2 avec doute ; le même terme, à la vue 230 et dans
la seconde accolade ci-dessus, est écrit $\frac{1}{\lambda}$. Une tache
d'encre est posée au milieu de la ligne à points de la seconde série.}

\marginal{D'après les valeurs de $\text{sin.}\,\omega$ rapportées §. 2 --}
En admettant que la condition $\text{sin.}\,\omega=0$ doivent s'accorder avec
toutes les solutions du problème on voit qu'il faut necessairement que
$\text{sin.}\,\lambda\omega=0$ ou que les facteurs
\[ \frac{1}{\lambda}-\frac{1-\lambda^2}{2.3.\lambda^3}\,\text{sin.}^2\lambda\omega+\text{\&c.},\qquad \Big\{\frac{1}{\lambda}-\frac{1-4\lambda^2}{2.3.\lambda^3}\,\text{sin.}^2\lambda\omega+\text{\&c}\Big\}\cos.\lambda\omega \]
suivant les cas, soient egaux a zéro. je suppose d'abord
$\text{sin.}\,\lambda\omega=0$, le facteur dont il s'agit se réduit a
$\frac{1}{\lambda}(1-e^{2(1-\lambda)\omega})$ car
$\text{sin.}(1-\lambda)\omega=\text{sin.}\,\omega\cos.\lambda\omega-\cos.\omega\,\text{sin.}\,\lambda\omega=-\text{sin.}\,\lambda\omega=0$,
et il est clair que $\frac{1}{\lambda}(1-e^{2(1-\lambda)\omega})$ ne peut
être supposé égal a zéro que dans la seule supposition $\omega=0$, mais si on
adopte cette dernière hypothèse le second facteur de l'équation du problème
se réduit a
\[ 2\big\{1+e^{2\lambda\omega}+e^{4\lambda\omega}+e^{6\lambda\omega}\ \text{\&c}+e^{2(1-\lambda)\omega}\big\}\,\text{sin.}(1-\lambda)\omega \]
\struck{qui} et il ne peut être par conséquent egale a zéro que dans la
supposition $\text{sin.}(1-\lambda)\omega=0$ supposition qui a cause de
$\text{sin.}\,\omega=0$ donne $\text{sin.}\,\lambda\omega=0$, ainsi il
necessairement admettre cette dernière.
\note{« et il » est écrit au-dessus de « qui » biffé, et « par conséquent »
au-dessus de la ligne entre « être » et « egale ». Dans la série qui suit
« se réduit a », une tache couvre un signe entre $e^{6\lambda\omega}$ et
$e^{2(1-\lambda)\omega}$ ; ce qui en dépasse se lit « \&c ». Un long trait
horizontal court sur « ainsi il necessairement admettre cette dernière » ; il
a la forme de la barre du $t$ rejetée et non d'une rature, et la phrase,
incomplète comme elle est écrite, est laissée.}

Je suppose donc $\text{sin.}\,\lambda\omega=0$ ; \struck{et} le facteur dont
il s'agit se réduit alors a
$\frac{1}{\lambda}(1-e^{2(1-\lambda)\omega})$ quantité qui ne peut être
egalée a zéro que dans la seule supposition $\omega=0$. J'ai déjà fait cette
remarque pur le cas ou $\lambda=\frac{1}{2}$ mais je vais prouver deplus
\marginal{$\times$ lorsqu'on a egard a la condition $\text{sin.}\,\omega=0$}
que dans le cas général, nonseulement $\times$ le second facteur de
l'équation du probleme ne peut être egalé a zéro que dans l'hypothèse
$\omega=0$, qui se rapporte au cas du repos de la lame, \struck{mais}
lorsqu'on a egard a la condition $\text{sin.}\,\omega=0$
\note{« pur » est écrit ainsi pour « pour ». La croix qui suit
« nonseulement » appelle la note portée dans la marge de gauche, sur une
ligne, qui s'ouvre elle-même sur la même croix ; elle est mise ici à sa
place, et la dernière ligne du feuillet la répète mot pour mot. « cas du »
est ajouté au-dessus de la ligne, entre « au » et « repos ». Le feuillet
s'arrête sur « $\text{sin.}\,\omega=0$ » ; le bas de la page est blanc.}

\page{238}

\note{au milieu du haut de la page, le nombre 10, souligné ; le nombre 116 est
à l'encre en haut à droite. Le verso, vue 239, est blanc.}

mais encore \struck{que sans avoir egard} qu'il y a des cas ou
qu'independamment de cette condition $\text{sin.}\,\omega=0$ l'équation
qu'il s'agit de resoudre n'est susceptible d'aucune solution qui puisse se
rapporter a cas $V$ c'est adire qui puisse mener a considerer les deux
portions dela lame separées par le stilet comme fixées a ce point et
appuiées a leur autre extremité $\times$
\marginal{$\times$ comme Euler le suppose \struck{dans} sa seconde solution.}
\note{« qu'il y a des cas ou » est écrit au-dessus des mots biffés. Après
« a cas », le $V$ est suivi d'un petit signe empâté qui n'est pas lu. La
croix qui ferme la phrase appelle la note portée dans la marge de gauche, en
deux lignes, qui s'ouvre elle-même sur la même croix ; elle est mise ici à
sa place.}

5 Pour cela je reprend l'équation finale
\[ 0=2(1-e^{2\omega})-(e^{\lambda\omega}-e^{(2-\lambda)\omega})(e^{\lambda\omega}-e^{-\lambda\omega})\frac{\text{sin.}\,\omega}{\text{sin.}\,\lambda\omega\,\text{sin.}(1-\lambda)\omega} \]
et je la transforme comme il suit :
\begin{align*}
0&=\frac{2(e^{\omega}-e^{-\omega})}{(e^{(\lambda-1)\omega}-e^{(1-\lambda)\omega})(e^{\lambda\omega}-e^{-\lambda\omega})}\cdot \text{sin.}\,\lambda\omega\,\text{sin.}(1-\lambda)\omega-\text{sin.}\,\omega \\
0&=\frac{(e^{(\lambda-1)\omega}+e^{(1-\lambda)\omega})(e^{\lambda\omega}-e^{-\lambda\omega})+(e^{\lambda\omega}+e^{-\lambda\omega})(e^{(1-\lambda)\omega}-e^{(\lambda-1)\omega})}{(e^{(\lambda-1)\omega}-e^{(1-\lambda)\omega})(e^{\lambda\omega}-e^{-\lambda\omega})}\,\text{sin.}\,\lambda\omega\,\text{sin.}(1-\lambda)\omega \\
&\qquad -\big\{\text{sin.}\,\lambda\omega\cos.(1-\lambda)\omega+\text{sin.}(1-\lambda)\omega\cos.\lambda\omega\big\} \\
0&=\frac{e^{(\lambda-1)\omega}+e^{(1-\lambda)\omega}}{e^{(\lambda-1)\omega}-e^{(1-\lambda)\omega}}\cdot\text{sin.}\,\lambda\omega\,\text{sin.}(1-\lambda)\omega-\cos.(1-\lambda)\omega\,\text{sin.}\,\lambda\omega \\
&\qquad +\frac{e^{\lambda\omega}+e^{-\lambda\omega}}{e^{\lambda\omega}-e^{-\lambda\omega}}\,\text{sin.}\,\lambda\omega\,\text{sin.}(1-\lambda)\omega-\text{sin.}(1-\lambda)\omega\cos.\lambda\omega
\end{align*}
et enfin
\begin{align*}
0&=\text{sin.}\,\lambda\omega\Big\{\frac{e^{(1-\lambda)\omega}+e^{(\lambda-1)\omega}}{e^{(1-\lambda)\omega}-e^{(\lambda-1)\omega}}\,\text{sin.}(1-\lambda)\omega-\cos.(1-\lambda)\omega\Big\} \\
&\qquad +\text{sin.}(1-\lambda)\omega\Big\{\frac{e^{\lambda\omega}+e^{-\lambda\omega}}{e^{\lambda\omega}-e^{-\lambda\omega}}\,\text{sin.}\,\lambda\omega-\cos.\lambda\omega\Big\}
\end{align*}
\note{le chiffre qui ouvre l'alinéa est la cinquième des marques de la
main ; il est tracé comme un $S$ et lu 5. Les deux signes « $-$ » posés
après $\text{sin.}(1-\lambda)\omega$ à la première transformation et devant
l'accolade à la seconde sont d'un trait très épais, presque une tache ; ils
sont rendus par un moins ordinaire.}

si on pouvoit considerer le stilet placé en $L$ comme fixé il faudroit
\struck{\uncertain{donc}} que les deux conditions
\[ \frac{e^{(1-\lambda)\omega}+e^{(\lambda-1)\omega}}{e^{(1-\lambda)\omega}-e^{(\lambda-1)\omega}}\,\text{sin.}(1-\lambda)\omega=\cos.(1-\lambda)\omega, \quad\text{et}\quad \frac{e^{\lambda\omega}+e^{-\lambda\omega}}{e^{\lambda\omega}-e^{-\lambda\omega}}\,\text{sin.}\,\lambda\omega=\struck{\uncertain{\text{cos.}}}\lambda\omega \]
ou
\[ \text{tang.}(1-\lambda)\omega=\frac{e^{(1-\lambda)\omega}-e^{(\lambda-1)\omega}}{e^{(1-\lambda)\omega}+e^{(\lambda-1)\omega}} \quad\text{et}\quad \text{tang.}\,\lambda\omega=\frac{e^{\lambda\omega}-e^{-\lambda\omega}}{e^{\lambda\omega}+e^{-\lambda\omega}} \]
eussent lieu en même tems, \struck{\ill{}} on auroit alors \struck{\ill{}}
une solution de l'équation qui \struck{\uncertain{seroit} \ill{} au même
tems} pourroit être rapportée au cas V. mais on voit aisément que ces deux
conditions nepeuvent \struck{avoir} exister a la fois dans certains cas ---
\struck{\uncertain{pour}} aucune autre supposition que celle $\omega=0$.
\note{dans la seconde des deux conditions, le mot biffé devant
$\lambda\omega$ est noyé sous la rature ; le calcul demande « cos. », et
c'est la lecture offerte, avec doute. Les autres ratures de cet alinéa ne se
laissent pas lire.}

En effet il faut aumoins que les deux portions de la lame rendent le même
son, car cette supposition est fondamentale

Or, si l'équation
$\text{tang.}\,\lambda\omega=\frac{e^{\lambda\omega}-e^{-\lambda\omega}}{e^{\lambda\omega}+e^{-\lambda\omega}}$
a lieu pour la portion $E L$ de longueur $\lambda a$ les sons simples que
cette portion \struck{\ill{}}
\note{le feuillet s'arrête sur ce mot biffé, qui n'est pas lu ; le bas de la
page est blanc.}

\page{240}

\note{au milieu du haut de la page, le nombre 11, souligné ; le nombre 117
est à l'encre en haut à droite. Une déchirure ancienne traverse le feuillet
en oblique, du haut du second alinéa au bas du troisième, sans emporter de
texte. Le verso, vue 241, n'appartient pas à ce lot.}

pourra rendre seront exprimés par les nombres suivans:
\[ \frac{25\pi c\surd 2gb}{16\lambda\lambda aa};\ \frac{81\pi c\surd 2gb}{16\lambda\lambda aa};\ \frac{169\pi c\surd 2gb}{16\lambda\lambda aa};\ \frac{289\pi c\surd 2gb}{16\lambda\lambda aa};\ \text{\&c.} \]
(voyez memoire d'Euler cas V)

et en supposant en même tems que l'autre équation
$\text{tang.}(1-\lambda)\omega=\frac{e^{(1-\lambda)\omega}-e^{(\lambda-1)\omega}}{e^{(1-\lambda)\omega}+e^{(\lambda-1)\omega}}$
ait pareillement lieu pour la portion $L E'$ de longueur $(1-\lambda)a$, les
sons simples que cette portion pourra rendre seront exprimés par les
nombres:
\[ \frac{25\pi c\surd 2gb}{16(1-\lambda)^2aa};\ \frac{81\pi c\surd 2gb}{16(1-\lambda)^2aa};\ \frac{169\pi c\surd 2gb}{16(1-\lambda)^2aa};\ \text{\&c.} \]

Je vais montrer a present \struck{que dans aucun autre cas que $\lambda=\frac{1}{2}$}
qu'il n'est pas généralement vrai que certains nombres \struck{un nombre} de
la première suite \struck{ne peut être egal a celui qui} puissent egaux a des
nombres appartenans \struck{lui correspond} dans la seconde \struck{en ayant
egard au nombre} \struck{des nœuds que doit contenir la plus longue portion
$L E$}
\note{l'alinéa est repris trois fois. « pas » est écrit dans la marge de
gauche ; « généralement vrai que certains nombres » est ajouté au-dessus de
la ligne, « puissent egaux a des nombres appartenans » au-dessous ; les mots
biffés sont donnés à leur place. La phrase ainsi recomposée reste
inachevée.}

soit pour plus de clarté $\lambda=\frac{1}{\lambda'}$ \quad
$(1-\lambda)^2=\frac{(\ill{}\lambda')^2}{\lambda'\lambda'}$ on voit
\struck{d'abord} alors facilement qu'aucun nombre de la première suite ne
peut être egal a un nombre quelconque dela seconde lorsque
$\lambda'-\uncertain{1}$ est un nombre pairement pair ou un nombre impair
dela forme $4n+3$ ainsi il ne reste \struck{donc a examiner} que les cas ou
$\lambda'-1$ est dela forme $4n+\uncertain{1}$, la plus petite valeur de
$\lambda'$ après 2 est 6 soit donc $\lambda'-1=5$, le plus petit nombre que
l'on puisse prendre dans la seconde suite est
$\frac{(25)^2(6)^2\pi c\surd 2gb}{16.5^2.aa}$ ce nombre est en effet egal au
premier de la première suite savoir a
$\frac{25(6)^2\pi c\surd 2gb}{16.aa}$
\struck{mais le nombre $25=4.6+1$ \uncertain{supposé} \ill{} dans les
portions} \struck{$L F$}
\note{« alors » et « facilement » sont écrits au-dessus de « d'abord »
biffé, les premières lettres de « facilement » repassées. Dans
$(1-\lambda)^2$, le début du numérateur est noyé sous une tache : le calcul
demanderait $\lambda'-1$, et rien n'est lu. Le chiffre après
« $\lambda'-$ » et celui qui ferme « $4n+$ » sont des taches pleines d'un
seul trait ; ils sont lus 1 avec doute, la suite du raisonnement les
appelant. La ligne biffée qui suit n'est lue qu'en partie.}

6 J'ai voulu entrer dans tous ces details afin d'être autorisée a conclure
qu'Euler a été induit en erreur lorsqu'il a cru pouvoir admettre sa seconde
solution dans laquelle il \struck{réduit} a supposer seulement que les deux
portions $E L$, $L F$ puissent se mouvoir en même tems sans exiger
\struck{qu'il \uncertain{soit} \ill{} \uncertain{souvenir} que la première
hypothese est que} \struck{le mouvement imprimé a} car \struck{cette
\uncertain{dernière}} j'ai \struck{\ill{}} \uncertain{voire} \struck{on}
voit que cette dernière condition même ne peut pas \struck{avoir lieu} être
remplie pour toutes les valeurs de $\lambda$.
\note{« seulement » est ajouté au-dessus de « supposer » ; « sans exiger »
au-dessus de la ligne biffée qui suit ; « j'ai \ill{} voire » au-dessus de
« on voit » biffé ; « être remplie » au-dessus de « avoir lieu » biffé. Les
deux lignes biffées de cet alinéa ne se laissent pas lire d'un bout à
l'autre. Le feuillet s'arrête sur « de $\lambda$. » ; le bas de la page est
blanc, et le lot s'arrête ici.}
\end{document}
